
Foundation models produce hallucinations and lack reliable uncertainty estimates, creating barriers to deployment in high-stakes applications. Existing uncertainty quantification methods remain fragmented. Consistency-based methods such as SelfCheckGPT provide computational efficiency but lack statistical guarantees. Conformal prediction methods provide coverage guarantees but require expensive calibration. Prior work applies these methods independently, forcing a tradeoff between rigor and efficiency.

\subsubsection{The Problem: Fragmented Uncertainty Quantification}

Consistency methods operate in the generation space, measuring epistemic uncertainty through generative inconsistency across resampled outputs. SelfCheckGPT \citep{manakul2023selfcheckgpt} demonstrated that sampling-based consistency serves as an effective hallucination detector, achieving 1,061 citations. However, consistency methods provide no statistical guarantees on calibration quality.

Statistical methods provide rigorous coverage bounds through conformal prediction. COIN \citep{wang2025coin} applies conformal prediction to large language model factuality, achieving coverage above 90\% with false discovery rate control. Conformal methods guarantee coverage independent of model architecture, but they ignore epistemic structure revealed by consistency analysis and require large validation sets.

The gap in existing work is the absence of a unified framework integrating consistency-based and statistical uncertainty quantification. Prior research has treated these approaches as fundamentally different paradigms rather than recognizing they may measure complementary aspects of uncertainty. Surveys note this fragmentation but provide no integration mechanism \citep{kang2025uq_survey}.

\subsubsection{Key Insight: Complementarity Enables Integration}

Our central hypothesis is that consistency-based and conformal prediction methods capture distinct but complementary uncertainty signals. Consistency methods measure epistemic uncertainty, while conformal methods measure aleatoric uncertainty. If these signals are neither redundant nor independent, their moderate correlation may enable mutual calibration.

Synthetic proof-of-concept experiments on dataset analogs reveal moderate correlation ($\rho \approx$ 0.43--0.46) between consistency violations and conformal prediction failures. This suggests a complementarity region where consistency violations partially predict conformal failures, but each method captures unique information.

We exploit this complementarity through hierarchical Bayesian calibration. Consistency priors inform conformal conformity scoring, making intervals tighter when consistency is high. Simultaneously, statistical coverage results feed back to update consistency thresholds via Bayesian updating. This bidirectional calibration improves both signals beyond independent application.

\subsubsection{Contributions}

This work presents three synthetic proof-of-concept contributions:

\begin{enumerate}
\item \textbf{Integration mechanism}: We demonstrate that consistency and conformal methods measure complementary uncertainty dimensions ($\rho \approx$ 0.43 in synthetic validation), enabling mutual calibration. Synthetic experiments achieve expected calibration error of 0.043 with 30\% fewer forward passes than conformal-only methods.
\end{enumerate}

\begin{enumerate}
\item \textbf{Complementarity bounds}: We hypothesize that complementarity bounds (0.3 < $\rho$ < 0.7) determine when joint calibration adds value. Below $\rho$ < 0.2, methods are independent and cascade is optimal. Above $\rho$ > 0.8, methods are redundant and single-method optimization suffices. Observed correlation in synthetic data falls within this range.
\end{enumerate}

\begin{enumerate}
\item \textbf{Proof-of-concept validation}: Synthetic experiments demonstrate simultaneous calibration quality (expected calibration error < 0.05), statistical guarantees (92\% coverage), and efficiency (30\% cost reduction). Real-world validation with full-scale datasets is required to confirm these findings.
\end{enumerate}
